\documentclass[a4paper,11pt]{article}
\usepackage[T1]{fontenc}
\usepackage[utf8]{inputenc}
\usepackage[slovene]{babel}
\usepackage[margin=2cm]{geometry}
\usepackage{amsmath,amssymb}
\usepackage{enumitem}
\usepackage{tabularx,array}
\usepackage{booktabs}
\usepackage{pgffor}
\usepackage{xcolor}
\usepackage{lmodern}
\usepackage{fancyhdr}
\usepackage{lastpage}

\setlength{\parindent}{0pt}
\setlength{\parskip}{4pt}
\setlist[enumerate]{itemsep=5pt,topsep=3pt}
\setlist[itemize]{itemsep=2pt,topsep=2pt}
\definecolor{modra}{HTML}{245B78}
\definecolor{barvaCrte}{HTML}{A6A6A6}
\pagestyle{fancy}
\fancyhf{}
\fancyfoot[C]{\small Vektorji \quad\textbullet\quad \thepage/\pageref{LastPage}}
\renewcommand{\headrulewidth}{0pt}
\newcommand{\tocke}[1]{\hfill\textit{#1/}}
\newcommand{\naloga}[2]{%
	\noindent\makebox[\textwidth]{%
		\parbox[t]{0.86\textwidth}{\textbf{#1}}% 
		\hfill\textit{#2/}%
	}\par
}
\newcommand{\odgovorcrte}[1]{%
	\par\vspace{1mm}%
	\begingroup
	\foreach\x in {#1}{%
		\noindent\makebox[\linewidth][l]{\rule{0pt}{0.62cm}\textcolor{barvaCrte}{\leaders\hrule height 0.35pt\hfill}}\par
	}
	\endgroup\vspace{1mm}%
}
\newcommand{\prostor}[1]{\par\vspace{#1}}

\begin{document}

\begin{center}
	{\LARGE\bfseries\color{modra}1.\ test: VEKTORJI}\\[3pt]
	\textbf{Čas pisanja: 50 minut}
\end{center}

\vspace{2mm}
\begin{tabularx}{\textwidth}{@{}X X@{}}
	\textbf{Ime in priimek:}\rule{0.2cm}{0pt}\rule{5cm}{0.4pt}&\textbf{Datum:}\rule{0.2cm}{0pt}\rule{3cm}{0.4pt}\tabularnewline
	\textbf{Letnik/oddelek:}\rule{0.2cm}{0pt}\rule{4.1cm}{0.4pt}&\textbf{Ocena:}\rule{0.2cm}{0pt}\rule{3.8cm}{0.4pt}
\end{tabularx}

\vspace{3mm}
\textbf{Točkovanje nalog:}
\begin{center}
\renewcommand{\arraystretch}{1.5}
\begin{tabular}{|>{\centering\arraybackslash}p{2.0cm}|*{7}{>{\centering\arraybackslash}p{1.2cm}|}>{\centering\arraybackslash}p{1.8cm}|}
\hline
\textbf{Naloga} & \textbf{1.} & \textbf{2.} & \textbf{3.} & \textbf{4.} & \textbf{5.} & \textbf{6.} & \textbf{7.} & \textbf{Skupaj} \\
\hline
\textbf{Možne točke} & 10 & 12 & 16 & 18 & 16 & 16 & 12 & 100 \\
\hline
\textbf{Dosežene točke} & & & & & & & & \\
\hline
\end{tabular}
\end{center}
\small\textit{Predlagana ocenjevalna lestvica: 0--49 točk: 1; 50--62: 2; 63--75: 3; 76--88: 4; 89--100: 5.}

\textbf{Navodila:} Piši čitljivo in prikaži postopek. Rezultate, kjer je smiselno, poenostavi. Pri geometrijskih nalogah lahko uporabiš skico; kalkulator ni potreben.

\naloga{1. naloga}{12}
V paralelogramu $ABCD$ označimo $\vec a=\vec{AB}$ in $\vec b=\vec{AD}$. Zapiši zahtevane vektorje s pomočjo $\vec a$ in/ali $\vec b$.
\begin{enumerate}[label=\alph*.]
	\item $\vec{AB}+\vec{BC}$ \hfill (3)
	\prostor{1.1cm}
	\item $\vec{AB}+\vec{AD}$ \hfill (3)
	\prostor{1.1cm}
	\item $\vec{AC}-\vec{AD}$ \hfill (3)
	\prostor{1.1cm}
	\item $\vec{BC}+\vec{CD}+\vec{DA}$ \hfill (3)
	\prostor{1.1cm}
\end{enumerate}

\newpage
\naloga{2. naloga}{16}
Dana sta vektorja $\vec a=(2,-1,3)$ in $\vec b=(-1,4,2)$.
\begin{enumerate}[label=\alph*.]
	\item Izračunaj $\vec a+\vec b$. \hfill (4)
	\prostor{1.2cm}
	\item Izračunaj $\vec a-\vec b$. \hfill (4)
	\prostor{1.2cm}
	\item Izračunaj $2\vec a-\dfrac12\vec b$. \hfill (4)
	\prostor{1.2cm}
	\item Točki sta $A(-1,2,0)$ in $B(3,-1,2)$. Določi komponente vektorja $\vec{AB}$. \hfill (4)
	\prostor{1.2cm}
\end{enumerate}


\newpage
\naloga{3. naloga}{18}
Dolžina, enotski vektor in skalarni produkt.
\begin{enumerate}[label=\alph*.]
	\item Izračunaj dolžino vektorja $\vec u=(3,4)$. \hfill (3)
	\prostor{1cm}
	\item Zapiši enotski vektor v smeri $\vec u$. \hfill (4)
	\prostor{1cm}
	\item Za $\vec p=(1,2,-2)$ in $\vec q=(2,-1,0)$ izračunaj skalarni produkt $\vec p\cdot\vec q$. \hfill (4)
	\prostor{1cm}
	\item Izračunaj $|\vec q|$. \hfill (3)
	\prostor{1cm}
	\item Utemelji, ali sta $\vec p$ in $\vec q$ pravokotna. \hfill (4)
	\odgovorcrte{1,2}
\end{enumerate}

\newpage
\naloga{4. naloga}{16}
Linearne kombinacije in vektorska enačba.
\begin{enumerate}[label=\alph*.]
	\item Dana sta $\vec u=2\vec a-\vec b$ in $\vec v=\vec a+3\vec b$. Izrazi $2\vec u-\vec v$ kot linearno kombinacijo vektorjev $\vec a$ in $\vec b$. \hfill (6)
	\prostor{1.8cm}
	\item Reši vektorsko enačbo $2(\vec x+\vec a)=3\vec b+\vec x$ in izrazi $\vec x$ z vektorjema $\vec a$ in $\vec b$. \hfill (6)
	\prostor{1.8cm}

\end{enumerate}

\newpage
\naloga{5. naloga}{16}
Odvisnost in baza. Podani so vektorji $\vec a=(1,2)$, $\vec b=(3,6)$ in $\vec c=(-2,1)$.
\begin{enumerate}[label=\alph*.]
	\item  Pokaži, da sta $\vec a$ in $\vec b$ linearno odvisna. \hfill (5)
	\odgovorcrte{1,2,3}
	\item Presodi, ali sta $\vec a$ in $\vec c$ linearno neodvisna. Odgovor utemelji. \hfill (4)
	\odgovorcrte{1,2,3}
	\item Kateri od parov $\{\vec a,\vec b\}$ in $\{\vec a,\vec c\}$ tvori bazo ravnine? Utemelji izbiro. \hfill (4)
	\odgovorcrte{1,2,3}
	\item Ali lahko trije vektorji v ravnini tvorijo bazo ravnine? Kratko pojasni. \hfill (3)
	\odgovorcrte{1,2}
\end{enumerate}

\vfill
\begin{center}
	\textit{Pred oddajo preveri, ali si odgovoril/-a na vse dele nalog.}
\end{center}

\end{document}
